\section{Computational methods}\label{app:methods}

This appendix records the algorithms behind the numerical statements of
\S\ref{sec:growth}, \S\ref{sec:monodromy} and \S\ref{sec:leaf}, the controls
that were run against them, and the residuals they achieved. It also discharges
the forward reference made in \S\ref{sec:monodromy}, where two algebraic
identities were deferred to here.

A statement of scope first, because it is easy to leave implicit and it should
not be. What follows states what was computed, at what precision, with which
controls, and to what residual. It does not make the computations
reproducible from this text alone: reproduction requires the scripts, which are
recorded with the claims and not reprinted here. That is the ordinary condition
of a paper of this kind, and it is said plainly rather than left to be
discovered.

\subsection{The solver, and its known-answer test}\label{ssec:methods-solver}

Every monodromy and Stokes computation in \S\ref{sec:monodromy} and
\S\ref{sec:leaf} is produced by one solver. Before it was applied to System
$\mathsf C$, it was run against a case whose answers are known in closed form.

\begin{remark}[Known-answer battery; \texttt{S22-01}, \textbf{VERIFIED}]
The solver reproduces Kummer's monodromy, formal and Stokes data against closed
forms: $32$ of $32$ checks pass at at least $35$ digits, with $44$ to $151$
digits observed, at both pre-registered parameter values.
\end{remark}

Two features of that test are worth naming. The parameter values were fixed
before the run, so the battery could not be tuned to the solver; and the
quantities checked include Stokes data, not only monodromy, so the irregular
part of the machinery is exercised rather than assumed.

\subsection{Truncation of the formal series}\label{ssec:methods-truncation}

The formal solution at the irregular point is divergent, and it is used
truncated. Where to truncate is decided by the series itself.

\begin{remark}[Optimal truncation; \texttt{S22-09}, \textbf{VERIFIED}]
The truncated formal series' residual sits at the minimal term of the series
--- ratio $1.29$ against a bound of $N^{3}$ --- and decays at the predicted
Gevrey rate $e^{-|t|R}$ between $R=160$ and $R=140$, with observed over
predicted $1.259$.
\end{remark}

The truncation radius is pre-registered rather than chosen after the fact, and
this has a consequence that surfaces in the transport computation: the
truncation level it fixes is a floor on the achievable agreement.

\begin{remark}[The floor is truncation, not precision; \texttt{S22-17}, \textbf{VERIFIED}]
Raising the working precision from $200$ to $320$ digits leaves the transport
agreement and its ratio unchanged to every printed figure --- $90.8118$,
$91.8301$, ratio $10.429999$ --- while every precision-limited control improves
by about $120$ digits: one residual from $1.20\times10^{-105}$ to
$5.57\times10^{-225}$, a second from $3.96\times10^{-112}$ to
$7.13\times10^{-232}$, a third from $2.35\times10^{-113}$ to
$3.30\times10^{-233}$, while the $\alpha$ residual stays at
$1.834\times10^{-91}$. The floor is the optimal-truncation level
$10^{-88.7}$ set by the pre-registered $R$, which no precision can reduce.
\end{remark}

This is the paper's clearest instance of a control that separates two
explanations of the same number. Had the residual been round-off, it would have
fallen with the precision; the controls did fall and it did not.

\subsection{Two algebraic checks, and what went wrong with them}\label{ssec:methods-algebra}

\S\ref{sec:monodromy} deferred two identities to this appendix. Both are
recorded here because each was, as first written, defective --- one vacuously
true, one comparing incompatible factorisation orders --- and because the
repaired versions do not carry the evidential weight the originals appeared to.

\begin{remark}[A check that could not fail; \texttt{S22-04}, \textbf{PROVEN}]
Check A6a as first written was algebraically vacuous: $(1+f_1f_2)/\alpha=M_{22}$
equals $(1-\det M)/m_{11}$ and so vanishes identically for any unimodular $M$
with $M_{11}\neq0$. It could not fail. This was proved by simplification to
zero together with an independent numeric arm.
\end{remark}

\begin{remark}[A check comparing two orderings; \texttt{S22-05}, \textbf{PROVEN}]
Check A7 as first written used the source's diagonal--lower--upper ordering,
with trace $\alpha+(1+f_1f_2)/\alpha$, against a computed
diagonal--upper--lower factorisation, with trace
$\alpha(1+f_1f_2)+1/\alpha$; the two differ by
$(\alpha-1/\alpha)f_1f_2$. At the known-answer values the wrong ordering misses
$2\cos(\pi b)$ by $0.9612$, while the corrected one holds to
$1.15\times10^{-29}$, reproducing the observed failure at order one.
\emph{The corrected identity is a dependent restatement of three checks
already made, not independent evidence.}
\end{remark}

The last sentence is the reason both appear here rather than among the
supporting evidence of \S\ref{sec:monodromy}. A repaired check that restates
what other checks already established is a consistency test on the algebra, and
counting it as a further confirmation would inflate the evidence.

\subsection{Extrapolation and high-precision limits}\label{ssec:methods-extrapolation}

The growth-law statements of \S\ref{sec:growth} rest on extrapolating a ratio
sequence, with the extrapolation itself controlled.

\begin{remark}[The exponent, over eight representatives; \texttt{S2-01}, \textbf{VERIFIED}]
For all $8$ quartic representatives the growth law
$Q_n/\Gamma(4n+1)\sim K(\beta_4/4^{4})^{n}n^{s}$ holds with
$s=\beta_3/\beta_4+3/2$. The maximum of
$\lvert s_{\text{measured}}-s_{\text{predicted}}\rvert$ over the eight is
$7.118\times10^{-16}$, after $4$ levels of ratio-$2$ Richardson extrapolation
on the ladder $n=1250,2500,5000,10000,20000$ at $80$ digits. The
pre-registered falsifier $\lvert\text{diff}\rvert>10^{-6}$ did not fire.
\end{remark}

\begin{remark}[The limit, to high precision; \texttt{S3B-09}, \textbf{VERIFIED}]
The limit of the normalised denominator sequence for the quadratic
representative was computed to $2317.8$ confirmed digits, by agreement between
runs at $N$ and $2N$ with $N=20000$, $L=640$, stride $520$, at $2800$ digits of
working precision. This supersedes an earlier $180$-digit value and an earlier
value at about $13.5$ digits; all three are mutually consistent.
\end{remark}

Two points of method are visible in the second of these. Confirmation is by
agreement between two runs at different truncation, not by a residual
internal to one run; and the superseded values are recorded as superseded and
checked for consistency, rather than replaced silently.

\subsection{What the methods do not establish}\label{ssec:methods-limits}

Three limitations follow from the above and are stated here rather than left
implicit.

\begin{enumerate}
\item Agreement between two runs at different truncation is evidence of
  internal consistency at a stated precision. It is not a proof, and every
  label attached to such a statement in this paper is \textbf{VERIFIED} for
  that reason.
\item The transport residual is bounded below by a pre-registered truncation
  radius, so the digits quoted in \S\ref{ssec:numerical-scope} are the
  instrument's ceiling at that step size, not an approach to an exact
  statement.
\item No displayed formula in this paper has been re-derived here from the
  script that computed it. The formulas trace to recorded claims; re-deriving
  them symbolically, and verifying against the page those that came from a page
  rather than a script, is a separate exercise this paper does not perform.
\end{enumerate}
